% !PRINTED 100
\weberpage{100}
%<<CONT
\[ \begin{aligned}
m = 819:\quad &\eta^3 - 273\eta + 91 = 0, & (\eta = \xi) \\
&\eta^3 - 273\eta + 91 . 19 = 0, \\
&\eta^3 - 273\eta + 91 . 8 = 0, \\
&\eta^3 - 273\eta + 91 . 17 = 0.
\end{aligned} \]

The roots of these individual equations are best found from equation (9)
or (11), by arranging them according to powers of $\varrho$, making use,
however, only of $\varrho^3 = 1$, not of $\varrho^2 + \varrho + 1 = 0$.

Thus one obtains, e.\,g., for the case $m = 63$:

\[ \begin{aligned}
(\varrho, \eta) &= (\eta_1 + \varrho\eta_1' + \varrho^2\eta_1'')(\eta_2 + \varrho\eta_2' + \varrho^2\eta_2'') \\
\text{or} \\
&= (\eta_1 + \varrho\eta_1' + \varrho^2\eta_1'')(\eta_2 + \varrho\eta_2'' + \varrho^2\eta_2'),
\end{aligned} \]

in which

\[ \begin{aligned}
\eta_1 &= r_1 + r_1^{-1}, & \eta_1' &= r_1^2 + r_1^{-2}, & \eta_1'' &= r_1^4 + r_1^{-4} \\
\eta_2 &= r_2 + r_2^{-1}, & \eta_2' &= r_2^3 + r_2^{-3}, & \eta_2'' &= r_2^2 + r_2^{-2}
\end{aligned} \]

are to be set (Vol.\,I, \S.\,172), and one thus finds:

\[ \begin{aligned}
\eta &= \eta_1\eta_2 + \eta_1'\eta_2'' + \eta_1''\eta_2' \\
\text{or} \\
\eta &= \eta_1\eta_2 + \eta_1'\eta_2' + \eta_1''\eta_2''.
\end{aligned} \]

In order to find the groups $\mathfrak{A}$ from these, one has to determine
two numbers $x, y$ in such a way that $7x + 9y = 1$, and then

\[ r = r_1 r_2, \quad \text{thus} \quad r_1 = r^{7x}, \ r_2 = r^{9y} \]

is to be set. If one takes $x = 4, y = -3$, one obtains for $x$ the two values % [sic: Weber prints ``für x'', although x has just been fixed and values of \eta follow.]
of $\eta$:

\[ \begin{aligned}
\eta &= (r^{28} + r^{-28})(r^{27} + r^{-27}) + (r^7 + r^{-7})(r^9 + r^{-9}) \\
&\qquad + (r^{14} + r^{-14})(r^{18} + r^{-18}) \\
\eta &= (r^{28} + r^{-28})(r^{27} + r^{-27}) + (r^7 + r^{-7})(r^{18} + r^{-18}) \\
&\qquad + (r^{14} + r^{-14})(r^9 + r^{-9}),
\end{aligned} \]

hence

\[ \mathfrak{A} = \pm 1, \ \pm 8, \ \pm 2, \ \pm 16, \ \pm 4, \ \pm 32, \]

or

\[ \mathfrak{A} = \pm 1, \ \pm 8, \ \pm 11, \ \pm 25, \ \pm 5, \ \pm 23. \]

The two groups can be represented by the powers of 2 and of 11.

The cubic equations whose laws of formation we have now come to know
contain, together with their \gesperrt{Tschirnhausen transformations, all
cubic cyclotomic equations.} But these equations are also all
%CONT>>
